%%%%%%%%%%%%%%%%%%%%%%%%%%%%%%%%%%%%%%
\section{Introduction}
%%%%%%%%%%%%%%%%%%%%%%%%%%%%%%%%%%%%%%
\lettrine{S}{ince} the creation of human-made orbital debris following the launch of Sputnik 1 in 1957, the first piece being the rocket stage that launched the artificial satellite and the second being the satellite itself, the presence of debris in orbit has increased dramatically.
Currently, about 27,000 officially cataloged objects, many of them non-functional and hence non-cooperative, are still in orbit and most of them are $10$ cm and larger. % https://www.nasa.gov/mission_pages/station/news/orbital_debris.html
Such non-cooperative systems in orbit lead to the threat of collision, potentially leading to the Kessler syndrome, the scenario where a sufficiently large collision of spacecraft could lead to a cascade effect and probably even make some low Earth orbits unusable for long-term use.
These collisions have involved operational satellites in the past such as the collision of Iridium 33, a communications satellite, with Cosmos 2251, an inactive Russian military communications satellite in 2009.
Our research is inspired by the need to control and deflect tumbling non-cooperative objects in orbit so as to avoid the risk of collisions with more vital orbiting systems such as functional satellites and space stations.

\cite{albee2021robust} presents techniques on rendezvous and control of such tumbling objects using a single spacecraft agent.
Although this has been conceptualized using single-agent systems, larger tumbling objects (hereon also referred to as 'targets') require multiple agents (hereon also referred to as 'chasers') to be brought to a standstill.
Furthermore, the paper assumes the accurate calculation of potentially unknown inertial and orbital parameters which may not always be true.
% In order to solve the problem of autonomous multi-agent interaction and deflection of uncooperative and adversarial objects in orbit, we consider @@@
As a result, we reconsider this problem considering the estimation of unknown parameters and the potential for extension to multiple chasers.

Previous research on on-orbit servicing has dealt with a similar problem where the chaser is attached to the target to perform repairs and similar operations \cite{nakka2022information,battistini2022iterative}.
However, such research assumes the cooperation of the target in terms of sharing information.
This is not the case with non-cooperative space debris that is inert and shut down in orbit.
We also draw motivation from previous literature on on-orbit inspection, especially using multiple chasers \cite{long2007orbit}.
However, this research delves into interaction with the target beyond inspection in order to stop the tumble of the target.

Chaser-target interaction can be made possible using several techniques.
These include the use of tethers by the chaser to attach onto a target that is at a distance \cite{zhang2019trajectory,olivieri2020test}.
While it is possible for an operational satellite or spacecraft to be either partially cooperative (via communication of target characteristics and dynamics) in the case of the on-orbit servicing problem, much of space debris is `non-cooperative and passive' \cite{albee2021robust}.
In addition, the tumble of the target would make it seemingly impossible to attach a free-flying tether at a specific location on the target.
Previous research \cite{park2017nonlinear} has dealt with this by instead proposing techniques to perform rendezvous and docking with the target.
Although this is a feasible solution, the possibility for collisions is very likely.
As a result, our research proposes rendezvous and docking along with the simultaneous estimation of unknown target parameters and the incorporation of robust techniques to prevent collisions.